\documentclass[11pt]{amsart} %


\usepackage{float}

\usepackage{epsfig}

\usepackage{array,amsmath, enumerate,  url, psfrag}
\usepackage{amssymb, amsaddr, fullpage}
\usepackage{graphicx,subfigure}
 \usepackage[usenames,dvipsnames]{pstricks}
\usepackage{color}

\restylefloat{figure}
%\usepackage{showkeys}
\renewcommand{\baselinestretch}{1.20}
%\topmargin =-0.4 in \oddsidemargin=0.25 in
\setlength{\textheight}{22.5 cm} \setlength{\textwidth}{16.2 cm}
\setlength{\unitlength}{5 mm} \makeatother



%\newcommand{\red}[1]{\textcolor{red}{#1}}
%\newcommand{\blue}[1]{\textcolor{blue}{#1}}

\newtheorem{theorem}{Theorem}[section]
\newtheorem{lemma}[theorem]{Lemma}
\newtheorem{conjecture}[theorem]{Conjecture}
\newtheorem{proposition}[theorem]{Proposition}
\newtheorem{corollary}[theorem]{Corollary}
\newtheorem{observation}[theorem]{Observation}


\newcommand{\n}{\noindent}
\newtheorem{thm1}{Theorem}
\newtheorem{satz}{thm1}

\newtheorem{thm2}{Theorem}
\newtheorem{sat}{thm2}

\newtheorem{thm}[sat]{Theorem}
\newtheorem{lem}[satz]{Lemma}
\newtheorem{cor}[satz]{Corollary}
\newtheorem{prop}[satz]{Proposition}
\newtheorem{ex}[satz]{Example}
\theoremstyle{definition}
\newtheorem{defn}[satz]{Definition}
\newtheorem{rem}[satz]{Remark}
\newtheorem{fic}{Figure}
\newtheorem{con}[satz]{Conjecture}
\newtheorem{clm}[satz]{Claim}

\newcommand{\di}{\displaystyle}


\iffalse

\newenvironment{definition}[1][Definition]{\begin{trivlist}
\item[\hskip \labelsep {\bfseries #1}]}{\end{trivlist}}
\newenvironment{example}[1][Example]{\begin{trivlist}
\item[\hskip \labelsep {\bfseries #1}]}{\end{trivlist}}
\newenvironment{remark}[1][Remark]{\begin{trivlist}
\item[\hskip \labelsep {\bfseries #1}]}{\end{trivlist}}

\fi

\title{Partitioning planar graphs without 4-cycles and 6-cycles 
into a forest and a disjoint union of paths}
\author{Pongpat Sittitrai$^{1}$ \hskip 0.2in Kittikorn Nakprasit$^{2}$}

\address{
	$^{1}$\small Department of Mathematics, Faculty of Science, Khon Kaen University, Khon Kaen, 40002, Thailand.\newline
	Email : pongpat.sittitrai@gmail.com\newline
	$^{2}$\small Department of Mathematics, Faculty of Science, Khon Kaen University, Khon Kaen, 40002, Thailand.\newline
	Email : kitnak@hotmail.com}
\bibliographystyle{plain}

\begin{document}

\maketitle

\begin{center}{\bf Abstract}\end{center}
\indent\indent

In this paper, we show that every planar graph without $4$-cycles 
and $6$-cycles has a partition of its vertex set into two sets, 
where one set induces a forest, and the other induces a forest 
with maximum degree at most $2$ 
(equivalently, a disjoint union of paths). 

Note that we can partition the vertex set of a forest into two independent sets. 
However a pair of independent sets combined may not induce a forest. 
Thus our result extends the result of Wang and Xu (2013) 
stating that the vertex set of every planar graph without $4$-cycles 
and $6$-cycles can be partitioned into three sets, 
where one induces a graph with maximum degree two, 
and the remaining two are independent sets. 


\section{Introduction} 
In this paper, we consider only undirected simple graphs. 
%A vertex partition is a widely studied topic in graph theory. 
Let $\mathcal{G}_i$ be a family of graphs. 
A graph $G$ with the vertex set $V(G)$ has a 
\emph{$(\mathcal{G}_1, \dots,  \mathcal{G}_k)$-partition} 
$(V_1, \dots, V_k)$ (a vertex partition) 
if $V(G)$ can be partitioned into 
$k$ sets $V_1$, $V_2,\dots, V_k$ where $V_i$ is an empty set or 
the induced subgraph $G[V_i]$ is in $\mathcal{G}_i$ for each $i\in\{1,\dots,k\}$. 

%In the rest of paper, we call a \emph{partition class $i$} as an induced subgraph of a vertex set $V_i$ in a  $(\mathcal{G}_1, \maththel{Gduced_2,\do $G[V_i]ts,   \mathcal{G}_n)$-partition of a graph. 

Certain classes of graphs are of interest. 
Let  $\mathcal{F}_d$ be a family of forests with maximum degree $d$, 
and let  $\Delta_d$ be a family of graphs with maximum degree $d$. 
We use $\mathcal{I}$ for $\mathcal{F}_0$ and $\Delta_0$, 
and we use $\mathcal{F}$ for $\mathcal{F}_\infty$ 
(a family of forests with unbounded degree). 
%Then, our paper consider only  $\mathcal{G}_i=\Delta_{d_i}$ or $\mathcal{F}_{d_i}$ for some $d_i$.  

Note that a $(\Delta_{d_1},\ldots,\Delta_{d_k})$-partition is 
equivalent to a $(d_1,d_2,\dots, d_n)$-coloring. 
Accordingly, an improper vertex coloring 
(a generalization of a proper coloring) 
can be regarded as a kind of vertex partition 
%Since every bipartite graph can be partitioned into one or two independent sets, a $\mathcal{G_1,\ldots,G_k}$-partition of $G$ where $\mathcal{G_i}$ contains only bipartite graphs for each $i\in \{1,\ldots,k\}$ implies $G$ is $2k$-colorable. 

% strengthen the results based on $(\Delta_{d_1},\Delta_{d_2},\dots, \Delta_{d_n})$-partition (or improper vertex coloring).\\
%\begin{enumerate}
%	\item The vertex set of any forests can be partitioned into two independent sets, but a graph induced by all vertices from two distinct independent sets  is not necessarily to be a forest. It means $(\mathcal{F}, \mathcal{G}_3,\dots,  \mathcal{G}_n)$-partition is a generalization of  $(\mathcal{I}, \mathcal{I}, \mathcal{G}_3,\dots,  \mathcal{G}_n)$-partition. 

	
%	\item $\mathcal{F}_d\subset \mathcal{G}_d$ for any integer $d\geq 2$. It means $(\mathcal{F}_{d_1}, \mathcal{G}_2,\dots,  \mathcal{G}_n)$-partition is a generalization of  $(\Delta_{d_1}, \mathcal{G}_2,\dots,  \mathcal{G}_n)$-partition for $d_1\geq2$. 
%\end{enumerate}
 
The following table shows some known results about the existence of 
particular vertex partitions for some classes of planar graphs. 
Some results may be redundant since a forest  
can be partitioned into one or two independent sets 
and $\mathcal{F}_d$ is a subclass of $\Delta_d.$ 
Nonetheless, we still put original results about 
$(\Delta_{d_1}, \Delta_{d_2},\dots,  \Delta_{d_n})$-partition 
in the table for a chronological reason.
 
 
\begin{center}
	\begin{tabular}{ |c|c|c| } 
		\hline
		Classes of Planar graphs  & $(\Delta_{d_1}, \dots,  \Delta_{d_n})$-partition & $(\mathcal{F}_{d_1}, \dots,  \mathcal{F}_{d_n})$-partition \\ 
	\hline
	Planar graphs& $(\mathcal{I},\mathcal{I},\mathcal{I},\mathcal{I})$ [FCT] & $(\mathcal{I},\mathcal{F},\mathcal{F})$ \cite{IFF} \\ 
		 & $(\Delta_2,\Delta_2,\Delta_2)$ \cite{222}& $(\mathcal{F}_2,\mathcal{F}_2,\mathcal{F}_2)$ \cite{F2F2F2} \\ 
	\hline
	Planar graphs with girth 4& $(\mathcal{I},\mathcal{I},\mathcal{I})$ \cite{000C3} & $(\mathcal{F}_5,\mathcal{F})$ \cite{F5FG4} \\ 
	\hline
	Planar graphs with girth 5& $(\Delta_3,\Delta_4)$\cite{34G5}  & $(\mathcal{F}_3,\mathcal{F}_3)$ (forbid adjacent $5$-cycles) \cite{F33G5} \\ 
	&  & $(\mathcal{I},\mathcal{F})$ \cite{IFG5} \\ 
	\hline
	Planar graphs with girth 6& $(\Delta_1,\Delta_4)$  \cite{02o04o14G8} & $(\mathcal{F}_1,\mathcal{F}_4)$ \cite{F1F4G6}  \\ 
	& $(\Delta_2,\Delta_2)$ \cite{22G6} & $(\mathcal{F}_2,\mathcal{F}_2)$ \cite{F2F2G6} \\ 
	\hline
	Planar graphs with girth 7& $(\mathcal{I},\Delta_4)$ \cite{02o04o14G8} & $(\mathcal{I},\mathcal{F}_5)$  \cite{IF235}\\ 
	\hline
	Planar graphs with girth 8& $(\mathcal{I},\Delta_2)$ \cite{02o04o14G8} & $(\mathcal{I},\mathcal{F}_3)$ \cite{IF235} \\ 
	\hline
	Planar graphs with girth 10& $(\mathcal{I},\mathcal{I},\mathcal{I})$ \cite{IF235} & $(\mathcal{I},\mathcal{F}_2)$ \cite{IF235}  \\ 
	\hline
	\end{tabular}
\end{center}
 
 
 
In \cite{giveva3}, Chartrand and Kronk gave an example of 
a planar graph without an $(\mathcal{F},\mathcal{F})$-partition. 
In view of this, finding the sufficient conditions for planar graphs 
to have an $(\mathcal{F},\mathcal{F})$-partition has become 
an interesting topic ever since. 

Every planar graphs with girth $4$ has an $(\mathcal{F},\mathcal{F})$-partition 
by being  $2$-degenerate. 
On the other hand, Montassier and  Ochem \cite{NomnC3} 
showed that for each $d_1$ and $d_2$, there exists 
a planar graph with girth $4$ having no 
$(\Delta_{d_1},\Delta_{d_2})$-partitions. 
Thus the result of  partitioning the vertex set of  
a planar graph with girth $4$ into two forests cannot be improved 
in terms of the maximum degrees of both forests. 
At best, one may find $d_1$ such that each planar graph 
with girth $4$ has an $(\mathcal{F}_{d_1},\mathcal{F})$-partition. 
Dross et al. \cite{F5FG4} verified this holds for $d_1=5$. 
However the case for $d_1 \leq 4$ is still open. 
In particular, the case $d_1=0$ if true (an $(\mathcal{I},\mathcal{F})$-partition) 
would imply the result by Gr\"otzsch \cite{000C3} 
(a $(\mathcal{I},\mathcal{I},\mathcal{I})$-partition). 

A planar graph without $4$- and $6$-cycles is shown to have 
an $(\mathcal{I},\mathcal{I},\mathcal{I})$-partition 
if it has no $8$-cycles by Wang and Chen \cite{000C4C6C8}  
or it has no $9$-cycles by Kang et al. \cite{000C4C6C9}. 
Liu and Yu \cite{1FC4C6C8C9} improved both results by showing 
that such graphs have an  $(\mathcal{I},\mathcal{F})$-partition. 

Sittitrai and Nakprasit \cite{35o44C4C5} proved that 
every planar graph without $4$-cycles and $5$-cycle has 
a $(\Delta_4,\Delta_4)$-partition and a $(\Delta_3,\Delta_5)$-partition. 
Later Cho et al. \cite{F34C45}  extended the result by proving that 
each such graph has an $(\mathcal{F}_3,\mathcal{F}_4)$-partition.
%which is stronger by adding no cycles condition to an induced graph of each partition class, and reducing the maximum degree of an induced graph of some partition class.
 
Wang and Xu \cite{200C4C6} showed that every planar graphs without 
$4$-cycles and $6$-cycles has a $(\Delta_2,\mathcal{I},\mathcal{I})$-partition. 
In this work, we improve their result in the following theorem.  
% ($\mathcal{F}_2,\mathcal{F}$)-partition, which is stronger by using above two concepts. 
 
% it is  stronger by  showing the partition class $1$ in $\mathcal{F}_2$ intend of $\Delta_2$,  and  the two partition classes $2$ and $3$ in edgeless graphs to  one partition class in $\mathcal{F}$.  
 
 \begin{theorem}\label{main}
 	Every planar graph without $4$-cycles and $6$-cycles has 
 	an ($\mathcal{F}_2,\mathcal{F}$)-partition.
 \end{theorem} 
 
On the other hand, in 2022 Kang et al. \cite{100C4C6} showed that 
every planar graph without $4$-cycles and $6$-cycles has 
an $(\mathcal{F}_1,\mathcal{I},\mathcal{I})$-partition 
(equivalently, a $(1,0,0)$-coloring). 
Inspired by two above results, we put forth the following conjecture.    
\begin{conjecture} 	Every planar graph without $4$-cycles and $6$-cycles 
has an $(\mathcal{F}_1,\mathcal{F})$-partition.
\end{conjecture}	


\section{Structures of a minimal counterexample}

Before we proceed to accrue results, 
some notation is required as follows. 
A \emph{$k$-vertex} (respectively, \emph{$k^+$-vertex} 
and  \emph{$k^-$-vertex})  is a vertex of degree $k$ 
(respectively, at least $k$ and at most $k.$) 
The same notation is applied for faces. 
A $k$-vertex $u$ is a \emph{$k$-neighbor} of $v$ if 
$u$ is adjacent to $v.$ 
The boundary walk of a face $f$ is denoted by $b(f).$ 
A vertex $v$ and a face $f$ are incident if $v$ is on $b(f).$ 
%and we write $f=v_1v_2\dots v_k$ if a walk $v_1v_2\dots v_k$ is the boundary of $f$.
If a vertex $v$ not on $b(f)$ but is adjacent to a $3$-vertex 
$u$ on $b(f),$  then we call $f$ a \emph{pendent face} 
of a vertex $v$ and  $v$ is a \emph{pendent neighbor} 
of $u$ (with respect to $f$).   
We use $n_i(v)$ to denote the number of incident $i$-faces 
of a vertex $v$, 
and use $m_i(v)$ to denote the number 
of pendent  $i$-faces of a vertex $v$.
%Likewise, $f$ is a \emph{pendent $3$-face} or a \emph{pendent $5$-face} of $v$ if $f$ is a $3$-face or a $5$-face, respectively.   

Given a $3$-vertex $u$ incident to a $3$-face or a $5$-face $f$, 
we call $u$ a \emph{terrible $3$-vertex} of $f$ if it has a 
pendent $4^-$-neighbor, otherwise we call $u$ 
a \emph{non-terrible $3$-vertex}. 
A $3$-face $f$ is a \emph{poor $3$-face} 
if $f$ is incident to two terrible $3$-vertices. 

Let $G$ be a minimal counterexample of Theorem~\ref{main}. 
That is $G$ does not has a ($\mathcal{F}_2,\mathcal{F}$)-partition, 
but each proper subgraph $G'$ of $G$ has 
an ($\mathcal{F}_2,\mathcal{F}$)-partition $(V_1,V_2).$  
%For all proofs below, we assume that $G'$ has an ($\mathcal{F}_2,\mathcal{F}$)-partition $(V_1,V_2),$ that is $V(G')$ is partitioned into  $V_1$ and $V_2$ where $G[V_1]$ inducing a forest with maximum degree two (a disjoint union of paths) and  $G[V_2]$ inducing a forest.  
For a vertex $v\in V_1,$  
a neighbor of $v$ in $V_1$ is a \emph{$V_1$-neighbor}, 
and we say $v$ is \emph{$V_1$-saturated} 
if $v$ has two $V_1$-neighbors. 
Some properties of $G$ are obtained as follows.

\begin{lemma}\label{atleast3} Each vertex in $G$ is a $3^+$-vertex.
\end{lemma}
\begin{proof} Let $u$ be a vertex in $G$. 
Suppose to the contrary that $u$ is a $2$-vertex. 
(The case that $u$ is a $1^-$-vertex is similar). 
Then $G-v$ has a ($\mathcal{F}_2,\mathcal{F}$)-partition $(V_1,V_2)$.  
Let $v$ and $w$ are two adjacent vertices of $u.$ 
	
If $v\in V_1$ or $w\in V_1$, 
then $G$ has a ($\mathcal{F}_2,\mathcal{F}$)-partition 
$(V_1,V_2\cup \{u\}),$ a contradiction. 
	
If $v, w\in V_2$, then $G$ has a ($\mathcal{F}_2,\mathcal{F}$)-partition $(V_1\cup \{u\},V_2),$ a contradiction.
\end{proof}

Since $G$ contains neither $4$- nor $6$-cycles, 
we have the following observation. 

\begin{observation}\label{face} Let $f$ be a face in $G$. 
	\begin{enumerate}
		\item [\rm (i)] A face $f$ is not a $4$-face.  
		\item [\rm (ii)] If $f$ is a $3$-face, 
		then $f$ is not adjacent to a $6^-$-face.
	\end{enumerate}
\end{observation}



\begin{lemma}\label{forbidvertex} 
    Let $v$ be a vertex of $G.$ 
	\begin{enumerate}
		\item [\rm (i)] If $v$ is a $3$-vertex, 
		        then $v$ is adjacent to a $5^+$-vertex. 
		\item [\rm (ii)] If $v$ is a $5$-vertex, 
		        then $v$ is not incident to a poor $3$-face. 
		\item [\rm (iii)] If $v$ is a $6$-vertex, 
		        then $v$ is incident to at most one poor $3$-face. 
	\end{enumerate}	
\end{lemma}

\begin{proof} 
%Let $v$ be a $k$-vertex with neighbors $v_1, \dots,v_k.$   
	
	\noindent (i)  Suppose to the contrary that 
	each neighbor of $v$ is a $4^-$-vertex. 
	Consider ($\mathcal{F}_2,\mathcal{F}$)-partition $(V_1,V_2)$ 
	of $G-v.$ 
	Since each neighbor of $v$ is a $3^-$-vertex in $G-v$, 
	we may assume that each of them is in $V_1$ 
	but is not $V_1$-saturated or is in $V_2.$    
	
	If two or three neighbors of $v$ are in $V_1,$ 
	then $G$ has a ($\mathcal{F}_2,\mathcal{F}$)-partition 
	$(V_1,V_2\cup \{v\}),$ 
	otherwise  $G$ has an ($\mathcal{F}_2,\mathcal{F}$)-partition $(V_1\cup \{v\},V_2)$ since the only neighbor of $v$ in $V_1$ 
	(if exists) is not $V_1$-saturated. 
	We obtain a contradiction for both cases.
	
    (ii)  Suppose to the contrary that $v$ is incident 
    to a poor $3$-face $f$ with $b(f)= vv_1v_2.$ 
    From the definition, $v_1$ and  $v_2$ are terrible 
    $3$-vertices with pendent $4^-$-neighbors, 
    say $v'_1$ and  $v'_2$, respectively. 
	
	Consider $G-\{v_1,v_2\}$ with an  ($\mathcal{F}_2,\mathcal{F}$)-partition $(V_1,V_2).$  
	
	Since $v$, $v'_1$, and $v'_2$ are $3^-$-vertices in $G-\{v_1,v_2\}$, 
	we may assume that each of them is in $V_1$ 
	but is not $V_1$-saturated or is in $V_2.$  
	
	-	Let $v\in V_1$. 
	
	If $v'_1$ or $v'_2$ is in $V_1$, 
	then $G$ has a ($\mathcal{F}_2,\mathcal{F}$)-partition
	$(V_1,V_2\cup\{v_1,v_2\}),$  a contradiction.
	
	If $v'_1$ and $v'_2$ are in $V_2$, 
	then $G$ has a ($\mathcal{F}_2,\mathcal{F}$)-partition 
	$(V_1\cup\{v_1\},V_2\cup\{v_2\})$ 
	since $v$ is not $V_1$-saturated, a contradiction.
	
	
	-	Let $v\in V_2$. 
	
	If $v'_1$ and $v'_2$ are in $V_2$, then $G$ has a ($\mathcal{F}_2,\mathcal{F}$)-partition $(V_1\cup\{v_1,v_2\},V_2),$  a contradiction.
	
	If $v'_1$ is in $V_1$ and $v'_2$ is in $V_2$, 
	then $G$ has a ($\mathcal{F}_2,\mathcal{F}$)-partition 
	$(V_1\cup\{v_2\},V_2\cup\{v_1\}),$ a contradiction.
	
	If $v'_1$ and $v'_2$ are in $V_1$, 
	then $G$ has a ($\mathcal{F}_2,\mathcal{F}$)-partition 
	$(V_1\cup\{v_1\},V_2\cup\{v_2\})$ 
	since $v'_1$ is not $V_1$-saturated, a contradiction. 
	
	(iii)  Suppose to the contrary that 
	$v$ is incident to two poor $3$-faces
	with boundary walks $vv_1v_2$ and $vv_3v_4.$  
	From the definition,  $v_i$ is a terrible $3$-vertices 
	with pendent $4^-$-neighbors, say $v'_i$ 
	where $i\in \{1,2,3,4\}$.   
	
	Consider $G'=G-\{v_1,v_2,v_3,v_4\}$ with 
	a ($\mathcal{F}_2,\mathcal{F}$)-partition $(V_1,V_2).$  
	Since each $v'_i$ is a $3^-$-vertex	in $G'$, 
	we may assume that $v'_i\in V_1$ but 
	$v'_i$ is not $V_1$-saturated or $v'_i\in V_2$  
	where $i\in \{1,2,3,4\}.$ 
	Moreover, since $v$ is a $2$-vertex in $G',$ 
	we may assume that $v\in V_1$ 
	but has no $V_1$-neighbors or $v\in V_2.$ 
	The table below shows that $G$ has 
	an ($\mathcal{F}_2,\mathcal{F}$)-partition for all cases, 
	a contradiction. 

\begin{center}
	\begin{tabular}{ |c|c|c|c| } 
		\hline
		$v$ & $v_1$ and $v_2$ &  $v_3$ and $v_4$  
		& $(\mathcal{F}_{2},  \mathcal{F})$-partition \\ 
	\hline
	$v\in V_1$ & $v'_1 \in V_1$ or $v'_2 \in V_1$ 
	&$v'_3 \in V_1$ or $v'_4 \in V_1$
	&$(V_1,V_2\cup\{v_1,v_2,v_3,v_4\})$\\ 
	\hline
	$v\in V_1$ & $v'_1 \in V_1$ or $v'_2 \in V_1$ 
	&$v'_3 \in V_2$ and $v'_4 \in V_2$
	&$(V_1\cup\{v_4\},V_2\cup\{v_1,v_2,v_3\})$\\ 
	\hline
	$v\in V_1$ & $v'_1 \in V_2$ and $v'_2 \in V_2$ 
	&$v'_3 \in V_2$ and $v'_4 \in V_2$
	&$(V_1\cup\{v_2,v_4\},V_2\cup\{v_1,v_3\})$\\
	\hline
	
	
	$v\in V_2$ & $v'_1 \in V_1$ or $v'_2 \in V_1$ 
	&$v'_3 \in V_1$ or $v'_4 \in V_1$
	&$(V_1\cup\{v_1,v_2,v_3,v_4\},V_2)$\\ 
	\hline
	$v\in V_2$ & $v'_1 \in V_1$ or $v'_2 \in V_1$ 
	&$v'_3 \in V_2$ and $v'_4 \in V_2$
	&$(V_1\cup\{v_1,v_2,v_3\},V_2\cup\{v_4\})$\\ 
	\hline
	$v\in V_2$ & $v'_1 \in V_2$ and $v'_2 \in V_2$ 
	&$v'_3 \in V_2$ and $v'_4 \in V_2$
	&$(V_1\cup\{v_1,v_2,v_3\},V_2\cup\{v_2,v_4\})$\\ 
	\hline
	\end{tabular}
\end{center}

\end{proof}


\begin{lemma}\label{number}
Each $k$-vertex $v$ has upper bounds on $m_i(v)$ and $n_i(v)$ 
as follows. 

	\begin{enumerate}
		\item [\rm (i)] $n_3(v)\leq\lfloor\frac{k}{2}\rfloor$. 
		
		\item [\rm (ii)]  
		$
		n_5(v) \leq
		\begin{cases}
		0,        & \text{if } m_3(v)+2n_3(v)= k,\\
		k,     & \text{if } m_3(v)+2n_3(v)=0,\\
		k-m_3(v)-2n_3(v)-1        & \text{,otherwise}.\\
		\end{cases}
		$
		
		
		\item [\rm (iii)]  $m_3(v)\leq k-2 n_3(v)$.
		
		\item [\rm (iv)] $m_5(v)\leq k-m_3(v)-2m_3(v)$.
	\end{enumerate}	
\end{lemma}	
\begin{proof} 
	Let $v_1, \ldots, v_k$ be neighbors of a $k$-vertex $v.$
	
	Let $A$ be the set of  $v_i$  where 
	$v_i$ is incident to an incident $3$-face of $v$. 
	
	Let $B$ be the set of  $v_i$ where 
	$v_i$ is pendent $3$-neighbor of $v$. 
	
	Let $C$ be the set of  $v_i$  where 
	$v_i$ is incident to an incident $5$-face of $v$. 
	
	Let $D$ be the set of  $v_i$  where 
	$v_i$ is a pendent $5$-neighbor of $v$.
	
	Recall Observation~\ref{face}(ii) that 
	a $3$-face is not adjacent to a $6^-$-face, 
	we have the following properties.
	
	(a) $A\cap B=\emptyset$ and  
	$(A\cup B)\cap (C\cup D)=\emptyset.$ 
	Consequently each $|C|$ and $|D|$  
	is not more than $k-|A|-|B|$. 
	
	(b) If $v_i\in A$, then $v_i$ is incident to exactly one incident $3$-face of $v$. 
	Consequently $n_3(v) \leq \frac{|A|}{2}.$ 
	
	(c) If $v_i\in C$, then  $v_{i-1}$ or $v_{i+1}$ is in $C$. 
	Consequently $n_5(v)= |C|-1$ for $1\leq |C| \leq k-1.$
	
	The Lemma \ref{number} follows (a), (b), and (c).
\end{proof}

\section{Proof of Theorem \ref{main}}

\indent Suppose $G$ is a minimal counterexample to the theorem. 
The discharging process is as follows. 
Let the initial charge of a vertex $v$ in $G$ be $\mu(v)=2d(v)-6,$ 
and let the initial charge of a face $f$ in $G$ be $\mu(f)=d(f)-6$. 
Using Euler's formula $|V(G)|-|E(G)|+|F(G)|=2$ 
and  the Handshaking lemma, we have
$$\displaystyle\sum_{v\in V(G)}\mu(v)
+\displaystyle\sum_{f\in F(G)}\mu(f)=-12.$$

\indent Now, we establish a new charge $\mu^*(x)$ 
for all $x\in V(G)\cup F(G)$ by transferring charge 
from one element to another and 
the summation of new charge $\mu^*(x)$ remains $-12.$  
If the final charge  $\mu^*(x)\geq 0$ for all $x\in V(G)\cup F(G)$, 
then we get a contradiction and the proof is completed.\\
\indent The discharging rules are as follows.\\
\noindent (R1) Let $v$ be a $3$-vertex. \\
- $v$ gives charge $\frac{1}{3}$ to each incident $3$-face. \\
\noindent (R2) Let $v$ be a $4$-vertex. \\
- $v$ gives charge $1$ to each incident $3$-face. \\
- $v$ gives charge $\frac{1}{2}$ to each incident $5$-face.\\
\noindent (R3) Let $v$ be a $5^+$-vertex. \\
- $v$ gives charge $\frac{5}{3}$ to each incident non poor $3$-face or $\frac{7}{3}$ to  each incident poor $3$-face. \\
- $v$ gives charge $\frac{2}{3}$ to each pendent $3$-face.\\
- $v$ gives charge $\frac{1}{4}$ to each pendent $5$-face.\\
- $v$ gives charge $\frac{1}{2}$ to each incident $5$-face.\\
\noindent (R4) Let $f$ be a $7^+$-face. \\
- $f$ gives charge $\frac{1}{6}$ to 
each incident $3$-vertex incident to a $3$-face.

\indent It remains to show that resulting $\mu^*(x)\geq 0$ 
for all $x\in V(G)\cup F(G)$. 
It is clear that $\mu^*(x)\geq 0$ 
when $x$ is a $3$-vertex not incident to any $3$-faces 
or $x$ is a $6$-face. 
%%%%%%%%%%%%%%%%%

\noindent\textit{\textbf{CASE 1:}} Let $f$ be a $3$-face.\\
\indent Note that $f$ is incident to a $5^+$-vertices 
if $f$ is incident to a terrible $3$-vertex 
by Lemma \ref{forbidvertex}(i). 
It follows that $f$ has at most two incident terrible $3$-vertices. 

Let $k$ denote the number of incident non-terrible 
$3$-vertices of $f.$ 

\indent  If $f$ has no incident terrible $3$-vertices,  
then $\mu^*(f)= \mu(f)+(3-k)\times1
+n\times\frac{1}{3}+k\times\frac{2}{3}=0$ 
by (R1), (R2), and (R3).

\indent  If $f$ has one incident terrible $3$-vertex,  
then $\mu^*(f)= \mu(f)+\frac{1}{3}+\frac{5}{3}+(1-k)\times1 
+k\times\frac{1}{3}+k\times\frac{2}{3}=0$ by (R1), (R2), and (R3). 

\indent  If $f$ has two incident terrible $3$-vertices,   
then $\mu^*(f)= \mu(f)+2\times\frac{1}{3}+\frac{7}{3}=0$ 
by (R1) and (R3).


\noindent\textit{\textbf{CASE 2:}} Let $f$ be a $5$-face.

\indent  If $f$ has no incident $4^+$-vertices,  
then $f$ has five non-terrible $3$-vertices 
by Lemma \ref{forbidvertex} (i). 
Thus $\mu^*(f)= \mu(f)+5\times\frac{1}{4}>0$ by (R3).

\indent  If $f$ has an incident $4^+$-vertex,  
then $f$ has at least  two non-terrible $3$-vertices 
by  Lemma \ref{forbidvertex} (i).
Thus $\mu^*(f)= \mu(f)+2\times\frac{1}{4}+\frac{1}{2}\geq 0$ by  (R2) and (R3).

\indent  If $f$ has at least two incident $4^+$-vertices, 
then $\mu^*(f)= \mu(f)+2\times\frac{1}{2}\geq 0$ by (R2) and (R3).

\noindent\textit{\textbf{CASE 3:}} Let $f$ be a $7^+$-face.

Note that a $7$-face is incident to at most six $3$-vertices 
incident to a $3$-face by Observation \ref{face}. 
By (R4), $\mu^*(f)= \mu(f)-6\times\frac{1}{6}\geq 0$ 
for a $7$-face $f$, 
and $\mu^*(f)= \mu(f)-d(v)\times\frac{1}{6}>0$ 
for a $8^+$-face $f.$ 

\noindent\textit{\textbf{CASE 4:}} 
Let $v$ be a $3$-vertex incident to a $3$-face. 

By Observation \ref{face},  
$v$ incident to one $3$-face and two $7^+$-faces. 
Thus $\mu^*(v)= \mu(v)-\frac{1}{3}+2\times\frac{1}{6}=0$ 
by  (R1) and (R4).


\noindent\textit{\textbf{CASE 5:}} Let $v$ be a $4$-vertex. 

By (R2),  $v$ loses charge $n_3(v)\times1+n_5(v)\times\frac{1}{2}$. 
By Lemmas \ref{number} (i) and (ii), we have the following cases. 

If $n_3(v)=0$, then $n_5(v)\leq4$. 
Thus $\mu^*(v)\geq \mu(v)-4\times\frac{1}{2}=0.$

If $n_3(v)=1$, then $n_5(v)\leq1$. 
Thus $\mu^*(v)\geq \mu(v)-1\times1-1\times\frac{1}{2}=0.$

If $n_3(v)=2$, then $n_5(v)=0$. 
Thus $\mu^*(v)= \mu(v)-2\times1=0.$ 

\noindent\textit{\textbf{CASE 6:}} Let $v$ be a $5^+$-vertex and $2\times n_3(v)+m_3(v)=0$. 

By Lemma \ref{number} (ii) and (iv), $n_5(v)\leq d(v)$ and $m_5(v)\leq d(v).$

Thus, 
$\mu^*(v) \geq \mu(v)-d(v)\times\frac{1}{2}-d(v)\times\frac{1}{4}
=2d(v)-6-d(v)\times \frac{3}{4}=\frac{5}{4}\times d(v)-6>0$ 
by (R3) and $d(v)\leq 5$.

\noindent\textit{\textbf{CASE 7:}} 
Let $v$ be a $5$-vertex with $2\times n_3(v)+m_3(v)>0$. 

By Lemma \ref{forbidvertex} (ii), 
$v$ is not incident  to a poor $3$-face. 
Then $v$ gives charge $\frac{5}{3}$ to each incident $3$-face. 
By (R3), we have\\
$$
\mu^*(v)=2d(v)-6-n_3(v)\times\frac{5}{3}+m_3(v)\times\frac{2}{3}+n_5(v)\times\frac{1}{2}+m_5(v)\times\frac{1}{4}.
$$

\textit{\textbf{CASE 7.1:}} 
Suppose $2\times n_3(v)+m_3(v)=d(v)=5$. 

By Lemma \ref{number}, we have $n_3(v)\leq2$ and  $m_3(v)= d(v)-2\times n_3(v),$ and $n_5(v)=m_5(v)=0.$
\begin{align*}
\mu^*(v)&=2d(v)-6-n_3(v)\times\frac{5}{3}-(d(v)-2\times n_3(v))\times\frac{2}{3}\\
&=\frac{4}{3}\times d(v)-6-n_3(v)\times\frac{1}{3}\\
&\geq\frac{4}{3}\times 5-6-2\times\frac{1}{3}\\
&= 0.
\end{align*} 



\textit{\textbf{CASE 7.2:}} 
Suppose $0<2\times n_3(v)+m_3(v)<d(v)=5$.

By Lemma \ref{number}, we have $n_3(v)\leq2$,  
$n_5(v)\leq d(v)-2n_3(v)-m_3(v)-1$, 
and $m_5(v)\leq d(v)-2\times n_3(v)-m_3(v).$

\begin{align*}
\mu^*(v)&\geq 2d(v)-6-n_3(v)\times\frac{5}{3}-m_3(v)\times\frac{2}{3}-(d(v)-2n_3(v)-m_3(v)-1)\times\frac{1}{2}\\
& -(d(v)-2n_3(v)-m_3(v))\times\frac{1}{4}\\
&=\frac{5}{4}\times d(v)-\frac{11}{2}-n_3(v)\times\frac{1}{6}+m_3(v)
\times\frac{1}{12}\\
&\geq \frac{5}{4}\times d(v)-\frac{11}{2}-n_3(v)\times\frac{1}{6}\\
&\geq \frac{5}{4}\times 5-\frac{11}{2}-2\times\frac{1}{6}\\
&> 0.
\end{align*}


\noindent\textit{\textbf{CASE 8:}} Let $v$ be a $6$-vertex 
with $2\times n_3(v)+m_3(v)>0$. 

By Lemma \ref{forbidvertex} (iii), 
$v$ is incident to at most one poor $3$-face. 
Note that it is enough to consider only 
the case $v$ containing exactly one poor $3$-face. 

By (R3),  we have   
\begin{align*}
\mu^*(v)&=2d(v)-6-\frac{7}{3}-(n_3(v)-1)\times\frac{5}{3}
-m_3(v)\times\frac{2}{3}-n_5(v)\times\frac{1}{2}
-m_5(v)\times\frac{1}{4}\\
&=-\frac{2}{3}+(2d(v)-6-n_3(v)\times\frac{5}{3}
-m_3(v)\times\frac{2}{3}-n_5(v)\times\frac{1}{2}
-m_5(v)\times\frac{1}{4}).
\end{align*}

\textit{\textbf{CASE 8.1:}} Suppose $2\times n_3(v)+m_3(v)=d(v)=6$. 

By Lemma \ref{number} (iii) and (iv), $n_5(v)=m_5(v)=0.$ 

It follows from \textit{\textbf{CASE 7.1}} that
\begin{align*}
\mu^*(v)&=-\frac{2}{3}+(\frac{4}{3}\times d(v)-6-n_3(v)\times\frac{1}{3})\\
&=\frac{4}{3}\times d(v)-\frac{20}{3}-n_3(v)\times\frac{1}{3}\\
&\geq\frac{4}{3}\times 6-\frac{20}{3}-3\times\frac{1}{3}\\
&>0.
\end{align*}


\textit{\textbf{CASE 8.2:}} Let $0<2\times n_3(v)+m_3(v)<d(v)=6$.

It follows from \textit{\textbf{CASE 7.2}} that
\begin{align*}
\mu^*(v)&\geq-\frac{2}{3}+(\frac{5}{4}\times d(v)
-\frac{11}{2}-n_3(v)\times\frac{1}{6})\\
&=\frac{5}{4}\times d(v)-\frac{37}{6}-n_3(v)\times\frac{1}{6}\\
&\geq\frac{5}{4}\times 6-\frac{37}{6}-3\times\frac{1}{6}\\
&>0.
\end{align*}

\noindent\textit{\textbf{CASE 9:}} Let $v$ be a $7^+$-vertex 
with $2\times n_3(v)+m_3(v)>0$. 

Then $v$ gives charge at most $\frac{7}{3}$ to each incident $3$-face. 

By R(3),  we have 
\begin{align*}
\mu^*(v)&=2d(v)-6-n_3(v)\times\frac{7}{3}-m_3(v)\times\frac{2}{3}
-n_5(v)\times\frac{1}{2}-m_5(v)\times\frac{1}{4}\\
&=-n_3(v)\times\frac{2}{3}+(2d(v)-6-n_3(v)\times\frac{5}{3}
-m_3(v)\times\frac{2}{3}-n_5(v)\times\frac{1}{2}
-m_5(v)\times\frac{1}{4}).
\end{align*}

\textit{\textbf{CASE 9.1:}} Suppose $2\times n_3(v)+m_3(v)=d(v)$.

By Lemma \ref{number} (iii) and (iv), we have $n_5(v)=m_5(v)=0.$ 

It follows from \textit{\textbf{CASE 7.1}} that
\begin{align*}
\mu^*(v)&=-n_3(v)\times\frac{2}{3}+(\frac{4}{3}\times d(v)-6-n_3(v)\times\frac{1}{3})\\
&=\frac{4}{3}\times d(v)-6-n_3(v)\times1.
\end{align*}

If $v$ is a $7$-vertex, then $d(v)=7$ and $n_3(v)\leq3$. 
Thus $\mu^*(v)\geq\frac{4}{3}\times 7-6-3>0.$

If $v$ is a $8^+$-vertex, then $d(v)\geq 8$ 
and $n_3(v)\leq\lfloor\frac{d(v)}{2}\rfloor\leq \frac{d(v)}{2}$. 
Thus $\mu^*(v)\geq\frac{4}{3}\times d(v)-6-\frac{d(v)}{2}>0.$

\textit{\textbf{CASE 9.2:}} Suppose $0<2\times n_3(v)+m_3(v)<d(v)$.

By Lemma \ref{number}, we have $n_3(v)\leq\frac{d(v)}{2}$,  
$n_5(v)\leq d(v)-2n_3(v)-m_3(v)-1$, 
and $m_5(v)\leq d(v)-2\times n_3(v)-m_3(v).$

It follows from \textit{\textbf{CASE 7.2}} that
\begin{align*}
\mu^*(v)&\geq-n_3(v)\times\frac{2}{3}
+(\frac{5}{4}\times d(v)-\frac{11}{2}-n_3(v)\times\frac{1}{6})\\
&=\frac{5}{4}\times d(v)-\frac{11}{2}-n_3(v)\times\frac{5}{6}\\
&\geq \frac{5}{4}\times d(v)
-\frac{11}{2}-\frac{d(v)}{2}\times\frac{5}{6}\\
&=\frac{5}{6}\times d(v)-\frac{11}{2}\\
&>0 \indent\text{ for each }d(v)\geq 7.
\end{align*}


\indent Finally,  it  follows from all cases 
that $\sum_{x\in V(G)\cup F(G)}\mu^*(x)>0$, a contradiction. 
This completes the proof.







%%%%%%%%%%%%%%%%%%


%\section{Acknowledgments}
\begin{thebibliography}{99}
	

	
%	\bibitem{1introgeq3}
%	G. Chartrand, H.V. Kronk, C.E. Wall, 
%	The point-arboricity of a graph, 
%	Israel J. Math. 6 (1968) 169--175.
\bibitem{IFF}
O. V. Borodin, 
A proof of Grünbaum's conjecture on the acyclic 
$5$‐colorability of planar graphs, 
Dokl. Akad. Nauk SSSR. 231 (1976) 18–-20.

\bibitem{IFG5}
O.V. Borodin, A.N. Glebov, 
On the partition of a planar graph of girth $5$ 
into an empty and an acyclic subgraph, 
Diskret. Anal. Issledovanie Oper. 8(4) (2001) 34--53 (in Russian).

\bibitem{02o04o14G8}
O.V. Borodin and A.V. Kostochka,
Defective 2-coloring of sparse graphs, 
J. Combin. Theory S. B 104 (2014) 72--80.
 
\bibitem{F2F2G6}
G. G. Chappell, J. Gimbel, and C. Hartman, 
Threshold for path colorings of planar graphs, 
Algorithms Combin. 26 (2005) 435--454.
 
\bibitem{giveva3}
G. Chartrand, H.V. Kronk, 
The point-arboricity of planar graphs, 
J. London Math. Soc. 44 (1969) 612--616.
 
\bibitem{F1F4G6}
M. Chen, A. Raspaud, W. Yu, 
An  ($F_1,F_4$)-partition of graphs 
with low genus and girth at least $6$,
J. Graph Theory 99(2) (2021) 186--206.
 

\bibitem{F34C45}
E-K. Cho, I. Choi, B. Park,
Partitioning planar graphs without $4$-cycles and $5$-cycles 
into bounded degree forests,
Discrete Math. 344(1) (2021) 112172.

\bibitem{34G5}
I. Choi, G. Yu, X. Zhang,
Planar graphs with girth at least $5$ are $(3,4)$-colorable,
Discrete Math. 342(12) (2019) 111577.

\bibitem{222}
L.J. Cowen, R.H. Cowen,  D.R. Woodall,
Defective colorings of graphs in surfaces:
partitions into subgraphs of bounded valency, 
J. Graph Theory 10(2) (1986) 187--195.

\bibitem{F5FG4}
F. Dross, M. Montassier, A. Pinlou, 
Partitioning a triangle‐free planar graph 
into a forest and a forest of bounded degree, 
Eur. J. Combin. 66 (2017), 81--94.

\bibitem{IF235}
F. Dross, M. Montassier, A. Pinlou, 
Partitioning sparse graphs into an independent set 
and a forest of bounded degree, 
Electron. J. Comb. 25(1) 2018 P1.45.

\bibitem{22G6}
F. Havet, J.-S. Sereni,
Improper choosability of graphs and maximum average degree, 
J. Graph Theory 52  (2006) 181--199.

\bibitem{000C3}
H. Gr\"otzsch, 
Zur theorie der diskreten gebilde, VII, 
Ein Dreifarbensatz für dreikreisfreie Netze auf der Kugel, 
Wiss. Z. MartinLuther-Universitat, Halle-Wittenberg, 
Math, Nat. Reihe 8 (1959) 109--120.
	 
\bibitem{100C4C6}
Y. Kang, L. Jin, P. Liu, Y. Wang, 
$(1,0,0)$-colorability of planar graphs without cycles 
of length $4$ or $6$,
Discrete Math. 345(4) (2022) 112758.
	
\bibitem{000C4C6C9}
Y. Kang, L. Jin, Y. Wang, 
The $3$-colorability of planar graphs without cycles of 
length $4$, $6$ and $9$, 
Discrete Math. 339 (2016) 299--307.
	
\bibitem{1FC4C6C8C9}
R. Liu, G. Yu,
Planar graphs without short even cycles are near-bipartite,
Discrete Appl. Math. 284 (2020) 626--630.

\bibitem{NomnC3}
M. Montassier, P. Ochem, 
Near-colorings: non-colorable graphs and NP-completeness, 
Electron. J. Comb. 22 (2015) 1--57.
	
\bibitem{F2F2F2}
K. Poh, 
On the linear vertex-arboricity of a plane graph, 
J. Graph Theory 14(1) (1990) 73--75.
	
\bibitem{35o44C4C5}
P. Sittitrai, K.Nakprasit, 
Defective $2$-colorings of planar graphs without 
$4$-cycles and $5$-cycles, 
Discrete Math. 341(8) (2018) 2142--2150. 
	
\bibitem{000C4C6C8}
W. Wang, M. Chen, 
Planar graphs without $4$, $6$, $8$-cycles are $3$-colorable,  
Sci. China A 50(11) (2007) 1552--1562.
	
\bibitem{F33G5}	Y. Wang, D. Huang, S. Finbow,
On the vertex partition of planar graphs into forests 
with bounded degree,
Applied Mathematics and Computation 374 (2020) 125032.
	
\bibitem{200C4C6}
Y. Wang, J. Xu,
Planar graphs with cycles of length neither $4$ nor $6$ 
are $(2,0,0)$-colorable, 
Inf. Proc. Lett. 113(18) (2013) 659--663.  

\end{thebibliography}
\end{document}